\section*{Chapter \ref{ch:iv1:timed-test-8} --- Timed Test 8}
\markboth{\protect\omTocMark{}{Answers: Timed Test 8}}{\protect\omTocMark{}{Answers: Timed Test 8}}

\begin{solution}{iq:iv1:timed-test-8:1}
The first is more likely, with $0.518$ against $0.491$.

Use the complement\index{complement}. The first is $1 - (\frac56)^4 = 1 - \frac{625}{1296} = 0.518$. The second is $1 - (\frac{35}{36})^{24} = 0.491$.

The attraction of the wrong answer is that they look the same. In both, the number of tries times the chance per try is $\frac23$, since $4 \times \frac16 = 24 \times \frac{1}{36}$. But this linear rule counts overlaps as if they did not occur, and the many small tries overlap more than the few large ones: the expected number of pairs of hits is about $0.21$ against $0.17$. An exact calculation confirms both, and a simulation agrees.

In the limit of many small tries, the chance tends to $1 - e^{-2/3} = 0.487$, so the second is already close to it.

\iqwrong{Saying that the two are equal, because both have an expected count of $\frac23$. An expected count is not a probability.}
\end{solution}

\begin{solution}{iq:iv1:timed-test-8:2}
The probability is $\frac{9}{64} \approx 0.141$.

Direct count: the closed walks have either two opposite pairs of the same axis, or one pair on each axis. Two pairs on one axis: $\binom42 = 6$ orders for each of 2 axes, so 12. One pair on each axis: $4! = 24$ orders. That makes 36 walks out of $4^4 = 256$.

A second route rotates the grid by $45^\circ$\index{rotation trick}. Then each step moves both new coordinates by $\pm1$, independently. Each coordinate is a separate walk, and each returns after 4 steps with chance $\binom42/16 = \frac38$. The square of this is $\frac{9}{64}$.

An enumeration of all 256 walks confirms it. The chance is lower than the $\frac38$ of one dimension, as there is more room to wander.

\iqwrong{Counting only the walks that step out twice and come straight back. Walks such as right, up, left, down also close, and the order of the moves matters.}
\end{solution}

\begin{solution}{iq:iv1:timed-test-8:3}
The market is pricing in $4\%$ growth. At $5\%$, the price is about 133.

The growing perpetuity gives $P = \frac{D}{r - g}$\index{discounted dividends}. So $100 = \frac{4}{0.08 - g}$, and $g = 4\%$. A numerical sum of the discounted dividends gives the same price of 100.

At $g = 5\%$, $P = \frac{4}{0.03} = 133.3$. A single point of growth raises the price by a third. This is why long-horizon assets are so sensitive to growth and rates.

Limit check: as $g$ tends to $r$, the price goes to infinity, and the formula fails at or above $r$. In that case, the sum diverges, which means that the assumption, not the market, is wrong.

\iqwrong{Treating the 8\% as dividend yield alone, and so answering zero growth. The yield is 4\%, and growth must supply the rest.}
\end{solution}

\begin{solution}{iq:iv1:timed-test-8:4}
The probability is $\frac{5}{36}$.

A difference of 5 means both a 1 and a 6 appear. Use inclusion--exclusion\index{inclusion--exclusion} over the complement.

There are 216 outcomes. Those with no 1 number $5^3 = 125$, and the same for no 6. Those with neither number $4^3 = 64$. So the good outcomes are $216 - 125 - 125 + 64 = 30$. And $\frac{30}{216} = \frac{5}{36}$.

A direct count checks it: the third die is a 1 or a 6, with 3 orders each, which gives 6, or a value from 2 to 5, with 6 orders each, which gives 24. The total is 30.

\iqwrong{Counting $3! \times 6$ for the third die in all cases, which gives 36. When the third die repeats a 1 or a 6, there are fewer orders.}
\end{solution}

\begin{solution}{iq:iv1:timed-test-8:5}
The probability is $\frac34$.

Given $X = x$, the chance that $Y < 1 - x$ is $1 - x$. Conditioning on $X < \frac12$ makes $X$ uniform on $[0, \frac12]$, with mean $\frac14$. So the answer is $1 - \frac14 = \frac34$\index{conditioning}.

Geometrically, the conditioning region is a rectangle of area $\frac12$. The part of it under the line $x + y = 1$ is a trapezium of area $\frac38$. The ratio is $\frac34$.

A simulation agrees. A limit check: conditioning on $X < a$ for small $a$ gives $1 - a/2$, which tends to 1, and for $a = 1$ gives $\frac12$, the unconditional chance.

\iqwrong{Giving $\frac12$, the unconditional value. Knowing that $X$ is small makes the sum smaller, so the chance must go up.}
\end{solution}

\begin{solution}{iq:iv1:timed-test-8:6}
The probability is $\frac{7}{18}$.

Condition on the type of the first tile\index{conditioning on the first draw}. A double has chance $\frac{7}{28} = \frac14$. It matches the 6 other tiles with that number. A non-double has chance $\frac34$, and it matches 6 others on each of its two numbers, so 12 tiles.

Out of the 27 tiles left,
\[
\frac14 \cdot \frac{6}{27} + \frac34 \cdot \frac{12}{27} = \frac{6 + 36}{108} = \frac{7}{18}.
\]
An enumeration of the $\binom{28}{2}$ pairs gives 147 matching pairs out of 378, which is the same value.

\iqwrong{Counting 12 matches for every tile, which gives $\frac49$. A double has only one number to offer.}
\end{solution}

\begin{solution}{iq:iv1:timed-test-8:7}
The expected number is $e \approx 2.718$.

Let $N$ be the number of draws. We have $N > n$ exactly when the first $n$ draws are in increasing order, which has chance $\frac{1}{n!}$, by symmetry among the $n!$ orderings\index{tail sum for expectation}. So
\[
\mathbb{E}N = \sum_{n \ge 0}P(N > n) = \sum_{n \ge 0}\frac{1}{n!} = e.
\]
A simulation returns $2.718$.

Check small cases: $P(N = 2) = \frac12$, a descent at once, and $P(N = 3) = \frac12 - \frac16 = \frac13$. These two alone contribute 2 to the mean, so a total just under 3 is sensible.

The answer does not depend on the distribution, provided that it is continuous, because only the ordering matters.

\iqwrong{Treating each later draw as having chance $\frac12$ of being a descent, which gives 3 by a geometric law. A long increasing run makes the last value high, and so the next descent likelier.}
\end{solution}

\begin{solution}{iq:iv1:timed-test-8:8}
With fair games, A's share is $\frac{11}{16}$. With $p = 0.6$, it is about $0.821$.

Play all the games that could be needed, even if the match would stop early. At most 4 games decide it\index{extended game}. A wins exactly when A wins at least 2 of the 4, since B then cannot reach 3.

With fair games, the chance is $1 - \frac{1 + 4}{16} = \frac{11}{16}$.

With $p = 0.6$, the chance that A wins 0 or 1 games is $0.4^4 + 4 \times 0.6 \times 0.4^3 = 0.0256 + 0.1536 = 0.1792$. So the share is $0.8208$. A direct recursion over the score gives both figures.

Limit check: if A needed no more games, the share would be 1, and if both needed the same number, it would be $\frac12$.

\iqwrong{Splitting by games still needed, in the ratio 3 to 2, which gives $\frac35$. That ignores how the remaining games can fall.}

\iqwhy{This is the root of pricing by a tree: the value now is the weighted value of what can still happen.}
\end{solution}

\begin{solution}{iq:iv1:timed-test-8:9}
In 3 dimensions, the rate is $\frac{\pi}{6} \approx 0.524$, with $\frac{6}{\pi} \approx 1.91$ draws. In 10 dimensions, the rate is about $0.0025$.

The acceptance rate is the ratio of volumes\index{rejection sampling}. The ball has volume $\frac43\pi$, and the cube has side 2, so volume 8. That gives $\frac{\pi}{6}$. The number of draws is geometric, with mean equal to the inverse rate.

In 10 dimensions, the cube has volume $2^{10} = 1024$. And $\pi^5 \approx 306$, so the ball has $\frac{306}{120} = 2.55$. The ratio is $\frac{2.55}{1024} = 0.0025$, or about 400 draws per point.

A simulation confirms both. So rejection from a cube fails in high dimension, since almost all the cube lies in its corners. Instead, draw a normal vector, normalise it, and scale the radius.

\iqwrong{Expecting the rate to settle near $\frac12$, because it is so in low dimension. It falls faster than exponentially in the dimension.}
\end{solution}

\begin{solution}{iq:iv1:timed-test-8:10}
The expected number is $\frac45$.

Take one king and the four aces. Among these five cards, the king is equally likely to be in any position\index{symmetry}, so it comes first with chance $\frac15$.

So each king comes before all the aces with chance $\frac15$. Indicators add over the 4 kings, and this gives $\frac45$.

A simulation returns $0.80$. A second route: the non-aces before the first ace number $\frac{48}{5}$ on average, and kings are 4 of the 48, so the mean is $\frac{48}{5} \cdot \frac{4}{48} = \frac45$.

Note that the other 44 cards in the deck do not matter, and the first route never used them.

\iqwrong{Using $\frac14$ per king, which gives 1. Each king races against 4 aces, so five cards are in the race, not four.}
\end{solution}

\begin{solution}{iq:iv1:timed-test-8:11}
The expected number of stops is $\frac{4651}{1296} \approx 3.59$, and the probability is $\frac{5}{54} \approx 0.093$.

A floor gets no stop with chance $(\frac56)^5 = \frac{3125}{7776}$. Summing indicators over the 6 floors\index{indicator variables},
\[
\mathbb{E} = 6\left(1 - \frac{3125}{7776}\right) = \frac{4651}{1296}.
\]

All different: $6 \cdot 5 \cdot 4 \cdot 3 \cdot 2 = 720$ of the $6^5 = 7776$ choices, or $\frac{5}{54}$.

An enumeration over all the 7776 choices gives the same. A limit check: with many floors, the stops tend to 5, one per person. With one floor, the expression gives exactly 1.

\iqwrong{Answering 5, one for each person, with no allowance for shared stops. Each floor must be looked at, not each person.}
\end{solution}

\begin{solution}{iq:iv1:timed-test-8:12}
The $R^2$ with both is $\frac13$, and with one alone it is $\frac14$.

Using standardised variables\index{normal equations}, the coefficients solve $b_1 + \frac12b_2 = \frac12$ and $\frac12b_1 + b_2 = \frac12$. By symmetry, $b_1 = b_2 = \frac13$.

And $R^2 = \sum b_ic_i = \frac13 \cdot \frac12 + \frac13 \cdot \frac12 = \frac13$.

With one alone, it is the squared correlation, $\frac14$. So the second regressor adds only $\frac{1}{12}$. A numerical solve agrees.

Limit checks: if the regressors were uncorrelated, the $R^2$ would be $\frac14 + \frac14 = \frac12$, and if they were perfectly correlated, it would stay $\frac14$. Our value lies between the two, as it should.

\iqwrong{Adding the squares, which gives $\frac12$. That only holds for uncorrelated regressors, as the overlap must not count twice.}
\end{solution}
